\section{vServer}\label{subsec:vserver}

In diesem Kapitel wird auf die Fortentwicklung des vServer eingegangen, auf dem die Collectiqo Sammlungsplattform seit Ende des zweiten Praxismoduls öffentlich zugänglich gemacht wird.
Neben strukturellen Anpassungen der Bereitstellung enthält dieses Kapitel auch Abschnitte zu den getroffenen Netzwerkanpassungen und der Weiterentwicklung der Docker-Umgebung.

\subsection{Bereitstellung}\label{subsubsec:vserver-bereitstellung}

Die Bereitstellung des vServer enthält Konfiguration zur Außenanbindung wie ein NGINX-Reverse-Proxy und Zertifikatsverwaltung mittels Certbot, was für die lokale Instanz nicht notwendig und störend ist.
Im Zuge der neuen Arbeitsweise mit git-Branches wurde die Struktur der Bereitstellung im Repository angepasst.
Die aktive Fortentwicklung des vServers erfolgt nun über einen separaten Branch \textit{deployment}, hier werden alle Änderungen vorgenommen und getestet, die für die Bereitstellung auf dem vServer notwendig sind.
Sowohl im main-Branch als auch im deployment-Branch wurde die Struktur der Bereitstellung angepasst, um eine klare Trennung zwischen der lokalen Entwicklung und der Bereitstellung auf dem vServer zu ermöglichen.
Für eine lokale Ausführung der Anwendung wurde ein separater Ordner \textit{deploy/local} angelegt, der die notwendige docker-compose.yaml für die lokale Entwicklung enthält.
Die Bereitstellung auf dem vServer erfolgt über den Ordner \textit{deploy/vServer}, hier enthalten ist die Docker-Compose-Datei inklusive Reverse-Proxy und Certbot-Zertifikatsverwaltung.

\subsection{Netzwerkanpassungen}\label{subsubsec:vserver-netzwerk}

Die Netzwerkanpassungen am vServer sind ein wichtiger Bestandteil des sicheren Betriebs der Collectiqo-Plattform.
Im Folgenden wird auf die Netzwerkseparierung, die automatische Zertifikatserneuerung und das Aktivieren von IPv6 eingegangen.

\subsubsection{Netzwerkseparierung}\label{subsubsec:vserver-netzwerk-separation}

Am Ende des zweiten Praxismoduls stellte sich mittels der Überprüfung der offenen Ports (Portscanning) des vServers heraus, dass mehr Ports offen waren als notwendig beziehungsweise gewünscht.
Neben den Ports 80 und 443 (http- und https-Übertragungsprotokoll), die für den NGINX-Reverse-Proxy genutzt werden, waren auch die Ports der Datenbank-Container (MySQL und MongoDB) offen erreichbar.
Dies stellt ein Sicherheitsrisiko dar, da unbefugte Zugriffe auf die Datenbanken möglich wären.
Die Datenbanken sollten nur über den Docker-Container, der das Front- und Backend der Collectiqo-Plattform bereitstellt, und nicht direkt über das Internet erreichbar sein.
Die Netzwerkanpassungen beinhalten die Trennung zwischen öffentlichem und privatem Netzwerk, um die Sicherheit des vServers zu erhöhen.

Es wurde ein öffentliches Netzwerk \textit{public-network} und ein privates Netzwerk \textit{private-network} in der Docker-Compose-Datei definiert.
Das \textit{public-network} ist ein Standard-Bridge-Netzwerk, das für die Kommunikation zwischen Containern und nach außen genutzt wird.
Das \textit{private-network} ist ebenfalls ein Bridge-Netzwerk, aber mit der Option \textit{"internal: true"}, wodurch der Zugriff von außen blockiert wird und nur interne Container-Kommunikation möglich ist.
So werden Dienste voneinander und vom externen Netzwerk gezielt abgeschottet.

Die node.js-Anwendung, die das Front- und Backend der Collectiqo-Plattform bereitstellt, ist im privaten Netzwerk angebunden, ebenso die Datenbank-Container (MySQL und MongoDB).
Der NGINX-Reverse-Proxy ist sowohl im öffentlichen Netzwerk als auch im privaten Netzwerk angebunden um Anfragen aus dem Internet entgegenzunehmen und an die node.js-Anwendung weiterzuleiten.
Der Certbot-Container ist rein im öffentlichen Netzwerk angebunden, da er nur für die Zertifikatsverwaltung zuständig ist und keine direkte Kommunikation mit der node.js-Anwendung benötigt.
Der Container läuft nur temporär, wenn ein neues Zertifikat ausgestellt oder ein bestehendes Zertifikat erneuert werden muss.


\subsubsection{Zertifikatserneuerung \& IPv6}\label{subsubsec:vserver-zertifikat-ipv6}

SSL-Zertifikate von Let's Encrypt haben eine Gültigkeit von 90 Tagen und müssen regelmäßig erneuert werden.
Die automatische Erneuerung der Zertifikate erfolgt mithilfe des Certbot-Containers, der in der Docker-Compose-Datei konfiguriert ist.
Es wurde ein CronJob eingerichtet, der jeden Sonntag gegen 8 Uhr morgens den Certbot-Container startet, um die Zertifikate zu erneuern.
Die Ausführung des Containers erfolgt mit dem Befehl \textit{docker compose --env-file .env.docker run --rm certbot renew --force-renewal}.
Der Befehlsparameter \textit{--force-renewal} sorgt dafür, dass die Zertifikate auch dann erneuert werden, wenn sie noch gültig sind.
Damit wird eine kontinuierliche Verfügbarkeit der SSL-Zertifikate sichergestellt, ohne dass manuelle Eingriffe erforderlich sind.
Die Rate-Limits von Let'S Encrypt erlauben fünf Zertifikatserneuerungen pro Woche, was ausreichend ist, um die automatische Erneuerung zu gewährleisten.
Der CronJob wurde so konfiguriert, dass ebenfalls eine E-Mail-Benachrichtigung an den Administrator gesendet wird.
Sollte die Erneuerung fehlschlagen, wird dies also rechtzeitig bemerkt und kann behoben werden.

Bislang war die Collectiqo-Plattform nur über IPv4-Adressen erreichbar.
Die Unterstützung für IPv6 wurde nun aktiviert, um die Erreichbarkeit der Plattform über das Internet zu verbessern und die Nutzung von IPv6-Adressen zu ermöglichen.
Hierzu wurde beim Provider ein AAAA-Record für die Domain collectiqo.de und collectiqo.com eingerichtet, der auf die IPv6-Adresse des vServers verweist.
Die Konfiguration des NGINX-Reverse-Proxys wurde entsprechend in der nginx.conf angepasst, um Anfragen über IPv6 zu unterstützen.
Diese wurde so erweitert, dass sie sowohl IPv4- als auch IPv6-Anfragen über beide Protokolle (http und https), also über die Ports 80 und 443, entgegennehmen kann.

Die Konfiguration des Nginx-Revers-Proxys basierend der Datei \textit{nginx/nginx.conf} sieht an den besagten Stellen (gekürzte Darstellung) wie folgt aus:
\vspace{1em}
\begin{lstlisting}[label={lst:lst-nginx-config}]
    server {
        listen 80 default_server; #IPv4 HTTP
        listen [::]:80 default_server; #IPv6 HTTP
        server_name collectiqo.de www.collectiqo.de collectiqo.com www.collectiqo.com;
    }

    server {
        listen 443 ssl; #IPv4 HTTPS
        listen [::]:443 ssl; #IPv6 HTTPS
        server_name collectiqo.de www.collectiqo.de collectiqo.com www.collectiqo.com;
    }
\end{lstlisting}
\vspace{1em}

\subsection{Weiterentwicklung Docker}\label{subsubsec:vserver-docker}

Healthchecks in Docker Compose (delayed start for node.js ehen DBs are ready and healthy)
Die Docker-Umgebung wurde weiterentwickelt, um die Stabilität und Zuverlässigkeit der Collectiqo-Plattform zu verbessern.
Die wichtigste Änderungen war die Implementierung von Healthchecks für die Docker-Container.
Zusätzlich wurde ein automatischer Start und Shutdown der Container in Abhängigkeit vom Host-System implementiert und eine E-Mail-Funktion für Benachrichtigungen hinzugefügt.

\subsubsection{Healthchecks}\label{subsubsec:vserver-docker-healthchecks}

Im Nachgang an die Abgabe des zweiten Praxismoduls wurde bemerkt, dass die node.js-Anwendung nicht immer ordnungsgemäß funktionierte.
Nach Sichtung der Logs stellte sich heraus, dass die Anwendung manchmal startete, bevor die Datenbank-Container (MySQL und MongoDB) bereit waren.
Dies führte dazu, dass die node.js-Anwendung nicht auf die Datenbanken zugreifen konnte und somit beispielsweise kein Login möglich war, welcher von der MySQL-Datenbank abhängig ist.
Manuelle Tests zeigten im Anschluss, dass die node.js-Anwendung ordnungsgemäß funktionierte, wenn die Datenbanken mit 30 Sekunden Vorlauf gestartet wurden.
Eine dauerhafte Lösung für dieses Problem, möglichst automatisiert, wurde gesucht.
Die Lösung bestand darin, Healthchecks für die Docker-Container zu implementieren.
Die Healthchecks ermöglichen es, den Zustand der Container zu überwachen und sicherzustellen, dass sie ordnungsgemäß basierend auf einer hinterlegten Testfunktion funktionieren.

Es wurden entsprechende Healthchecks für beide Datenbank-Container (MySQL und MongoDB) sowie für die node.js-Anwendung und den NGINX-Reverse-Proxy in der Docker-Compose-Datei \textit{docker-compose.yaml} (gekürzte Darstellung) definiert:
implementiert.
\vspace{1em}
\begin{lstlisting}[label={lst:lst-dockercompose-config1}]
services:
  mysql:
    image: lorackdev/clq-mysql-multiarch
    restart: always
    healthcheck:
      test: ["CMD-SHELL", "mysqladmin ping -h localhost -u root -p${MYSQL_DATABASE_PASSWORD}"]
      start_interval: 30s
      start_period: 45s
      interval: 10s
      timeout: 5s
      retries: 5

  mongodb:
    image: mongo:latest
    restart: always
    healthcheck:
      test: ["CMD-SHELL", "mongosh --host localhost --port 27017 -u ${MONGO_DATABASE_USER} -p ${MONGO_DATABASE_PASSWORD} --eval 'db.runCommand(\"ping\").ok' --quiet"]

  web:
    image: clq-app
    restart: always
    healthcheck:
      test: ["CMD-SHELL", "ps aux | grep node | grep -v grep || exit 1"]
    depends_on:
      mysql:
        condition: service_healthy
      mongodb:
        condition: service_healthy

  nginx-reverse-proxy:
    image: nginx:latest
    restart: unless-stopped
    healthcheck:
      test: ["CMD-SHELL", "curl -sf http://localhost/health || exit 1"]
    depends_on:
      web:
        condition: service_healthy
\end{lstlisting}
\vspace{1em}

Dabei wird bei der \textbf{MySQL} Datenbank mit dem Befehl \textit{mysqladmin ping} geprüft, ob der MySQL-Server auf localhost mit dem Benutzer root und dem Passwort aus der Umgebungsvariable MYSQL\_DATABASE\_PASSWORD erreichbar ist.
Liefert der Befehl keine erfolgreiche Antwort, gilt der Container als nicht gesund.

Für die \textbf{MongoDB} Datenbank wird mit dem Befehl \textit{mongosh} abgefragt, ob eine Verbindung zur lokalen MongoDB-Instanz auf Port 27017 möglich ist und ob der Datenbankbefehl \textit{db.runCommand("ping").ok} erfolgreich ausgeführt werden kann.
Ist das der Fall, gilt der Container als „healthy“, andernfalls wird ein Fehlerstatus zurückgegeben.
Die Zugangsdaten werden dabei ebenfalls aus den Umgebungsvariablen entnommen.

Ob ein \textbf{Node.js}-Prozess im Container läuft wird mit dem Kommando \textit{ps aux | grep node | grep -v grep || exit 1} überprüft.
Wird keiner gefunden, liefert der Healthcheck einen Fehlerstatus zurück.
So erkennt Docker, ob der Service noch funktionsfähig ist.

Für den \textbf{NGINX} Container wird per curl erprobt, ob der HTTP-Endpunkt http://localhost/health erreichbar ist.
Schlägt der Aufruf fehl, liefert der Healthcheck einen Fehlerstatus zurück.
So kann Docker erkennen, ob der NGINX-Container ordnungsgemäß läuft.

Die Healthchecks in Docker Compose ermöglichen es, den Zustand der Container zu überwachen und sicherzustellen, dass sie ordnungsgemäß funktionieren.
Dabei kommen verschiedene Parameter zum Einsatz, welche bei allen Container gleich gewählt wurden.
\begin{itemize}
    \item start\_interval: Zeit zwischen den ersten Healthcheck-Versuchen nach Containerstart.
    \item start\_period: Zeitraum nach Start, in dem Fehler beim Healthcheck ignoriert werden (Wartezeit bis zur ersten Bewertung).
    \item interval: Zeit zwischen den einzelnen Healthcheck-Versuchen im Normalbetrieb.
    \item timeout: Maximale Zeit, die ein Healthcheck-Befehl laufen darf, bevor er als fehlgeschlagen gilt.
    \item retries: Anzahl der erlaubten Fehlschläge, bevor der Container als „unhealthy“ markiert wird.
\end{itemize}

Mit der Option \textit{depends\_on} wird sichergestellt, dass die node.js-Anwendung erst gestartet wird, wenn beide Datenbank-Container als „healthy“ markiert sind.
NGINX wiederum startet erst, wenn die node.js-Anwendung als „healthy“ markiert ist.
Da Certbot ebenfalls eine Abhängigkeit zum NGINX-Reverse-Proxy hat, wird dieser ebenfalls erst gestartet, wenn der NGINX-Container als „healthy“ markiert ist.
Somit erfolgt eine verzögerte Startreihenfolge der Container, die sicherstellt, dass die Container in der richtigen Reihenfolge und Abhängigkeiten gestartet werden.

\subsubsection{Automatischer Start und Shutdown}\label{subsubsec:vserver-docker-auto-start}

Eine weitere Verbesserung der Docker-Umgebung ist die Implementierung eines automatischen Starts und Shutdowns der Container in Abhängigkeit vom Host-System.
Nach einem Neustart des Host-Systems war auffällig geworden, dass grade der NGINX-Reverse-Proxy in einem Restart-Zustand hängen blieb und nicht ordnungsgemäß startete.
Es wurde ein Systemd-Service  namens \textit{docker-compose-app.service} unter \textit{/etc/systemd/system/} erstellt, der die Docker-Container beim Hochfahren des Host-Systems automatisch startet und beim Herunterfahren des Host-Systems automatisch stoppt.
Die Systemd-Service-Datei sieht wie folgt aus:
\vspace{1em}
\begin{lstlisting}[label={lst:lst-systemd-service}]
[Unit]
Description=Docker Compose Application
Requires=docker.service
After=docker.service

[Service]
Type=oneshot
RemainAfterExit=yes
WorkingDirectory=/home/collectiqo/git/collectiqo/deploy/vServer
ExecStart=/usr/bin/docker compose --env-file .env.docker up -d
ExecStop=/usr/bin/docker compose --env-file .env.docker down --timeout 15
TimeoutStopSec=90
KillMode=process

[Install]
WantedBy=multi-user.target
\end{lstlisting}
\vspace{1em}

Innerhalb der Service-Datei wird das WorkingDirectory auf den Ordner gesetzt, in dem sich die Docker-Compose-Datei samt env-Datei befindet.
Die Befehle \textit{ExecStart} und \textit{ExecStop} starten und stoppen die Docker-Container in Kombination mit dem Host-System.
Ein \textit{TimeoutStopSec} von 90 Sekunden wurde gesetzt, um sicherzustellen, dass die Container genügend Zeit haben, um ordnungsgemäß herunterzufahren, bevor der Host-Computer heruntergefahren wird.
Die Option \textit{RemainAfterExit=yes} sorgt dafür, dass der Service als aktiv bleibt, auch wenn die Docker-Container gestoppt wurden.
Die Systemd-Service-Datei wurde abschließend aktiviert, sodass sie beim Booten des Host-Systems automatisch gestartet wird.


\subsubsection{E-Mail-Funktion}\label{subsubsec:vserver-docker-email}

In diesem Praxismodul wurde ein weiterer Container hinzugefügt, der eine Funktionalität für E-Mail bereitstellt.
Dieser Container fungiert als SendOnly-SMTP Mail Transport Agent (MTA) und ermöglicht es, E-Mails rein zu versenden, der Empfang von E-Mails ist bewusst nicht vorgesehen.
Der E-Mail-Container ist in der Docker-Compose-Datei \textit{docker-compose.yaml} definiert und nutzt das Image \textit{boky/postfix}, welches einen leichtgewichtigen Postfix-Mailserver bereitstellt.

Wie der NGINX-Reverse-Proxy ist der Mail-Container sowohl im öffentlichen Netzwerk als auch im privaten Netzwerk angebunden, andere Container beziehungsweise das Host-System können somit über den definierten Port 2525 E-Mails versenden.
Er stellt die Grundlage für eine mögliche Implementierung von E-Mail-Funktionen in der Collectiqo-Plattform dar, beispielsweise um ein Passwort-Zurücksetzung oder eine Authentifizierung über einen zweiten Faktor zu ermöglichen.


